\documentclass[a4paper,12pt]{article}
\usepackage[utf8]{inputenc}
\usepackage{hyperref}
\usepackage{geometry}
\geometry{margin=1in}

\title{mighf System Specification \\ \large Architecture, Operation, and Porting Guide}
\author{Miguel “Kuyo” / SufremOak}
\date{\today}

\begin{document}

\maketitle

\begin{abstract}
This document specifies the architecture and operational principles of the mighf system, including the instruction set, memory model, execution pipeline, and guidelines for reimplementation and porting to new hardware platforms.
\end{abstract}

\section{Overview}
The mighf system is a minimalist, stack- and register-based CPU architecture designed for easy implementation and efficient execution on embedded and legacy hardware. It balances simplicity and practical features while supporting extensibility and portability.

\section{Architecture Components}
\subsection{Registers}
\begin{itemize}
    \item \textbf{R0--R7}: 8 general-purpose 16-bit registers for computation and temporary storage.
    \item \textbf{PC} (Program Counter): Holds the address of the next instruction to execute.
    \item \textbf{SP} (Stack Pointer): Points to the top of the stack in memory.
    \item \textbf{FLAGS}: Status register containing Zero (Z), Carry (C), Sign (S), and Overflow (O) flags.
\end{itemize}

\subsection{Memory Model}
\begin{itemize}
    \item Flat 16-bit address space, byte-addressable.
    \item Stack grows downward from a high memory address.
    \item Supports direct memory access for load/store instructions.
\end{itemize}

\section{Instruction Set}
mighf uses fixed 16-bit wide instructions with opcode and operands encoded compactly. Supports arithmetic, logic, control flow, stack, and memory operations.

\section{Execution Model}
\begin{itemize}
    \item \textbf{Fetch}: CPU fetches the instruction pointed by PC.
    \item \textbf{Decode}: Instruction is decoded into opcode and operands.
    \item \textbf{Execute}: ALU and other units perform the operation.
    \item \textbf{Memory Access}: For load/store, memory is accessed as needed.
    \item \textbf{Write-back}: Results are stored back in registers or memory.
    \item \textbf{Update PC}: PC is incremented or updated by control flow instructions.
\end{itemize}

\section{System Behavior}
\begin{itemize}
    \item On reset, PC is set to a predefined start address.
    \item Interrupts and exceptions are outside the current ISA scope but can be added as extensions.
    \item The stack pointer must be initialized before use.
\end{itemize}

\section{Porting and Reimplementation Guidelines}
\subsection{Target Platform Requirements}
\begin{itemize}
    \item Minimum 16-bit addressable memory.
    \item Ability to implement 16-bit registers and ALU operations.
    \item Support for a basic stack mechanism.
\end{itemize}

\subsection{Reimplementation Strategies}
\begin{itemize}
    \item \textbf{Software Emulation}: Implement the ISA in a high-level language (C, Rust, Go) using a loop for fetch-decode-execute.
    \item \textbf{FPGA/Hardware}: Design the datapath, registers, and control logic following the ISA spec; leverage simple microcoded or hardwired control units.
    \item \textbf{Cross-Platform Considerations}:  
      \begin{itemize}
        \item Endianness: Define the instruction encoding byte order explicitly (big or little endian).
        \item Memory alignment: Enforce 16-bit aligned instructions and data where possible.
        \item Timing: Pipeline stages can be optimized or simplified depending on hardware.
      \end{itemize}
\end{itemize}

\subsection{Extensibility}
The architecture is designed for easy extension by adding:
\begin{itemize}
    \item New instructions with reserved opcode space.
    \item Expanded register sets or wider registers.
    \item Interrupt and exception handling.
    \item Coprocessor or peripheral integration.
\end{itemize}

\section{Example: Simple Reimplementation Loop (Pseudo C)}
\begin{verbatim}
while (running) {
    instr = fetch(memory, CPU.PC);
    decoded = decode(instr);
    execute(decoded, &CPU, memory);
    updatePC(&CPU, decoded);
}
\end{verbatim}

\section{Summary}
The mighf system is built to be straightforward to implement and maintain, with a clear focus on portability and extensibility. Its simple core makes it ideal for educational use, embedded systems, and experimental CPU projects.

\end{document}
